\noindent\textbf{BEMERKUNGEN. ZUR NETZAUSGLEICHUNG.}\hfill327

\begin{center}\textbf{BEMERKUNGEN.}\end{center}

\par\noindent\hspace*{\parindent}Von der Netzausgleichung sind nur einige lose, nicht nummerierte Blätter
vorhanden. Eine vollständige Rekonstruktion bis in alle Einzelheiten anhand
dieser Blätter --- sofern sie überhaupt möglich ist --- wäre mit
aufwendigen Rechnungen verbunden.

\par\noindent\hspace*{\parindent}Direkt aus diesen Blättern übernommen sind Art.~[2], die Normalgleichungen
des Art.~[4], die Seitengleichungen in Art.~[5] und die Tabelle unter [9],
der nur die Spaltenbezeichnungen hinzugefügt wurden. Nach diesen Blättern
zusammengestellt sind Art.~[3], die Tabellen der Art.~[7] und [8] sowie die
Normalgleichungen des Art.~[6]. Die Notiz [1] stammt aus einem Beobachtungs-
und Rechnungsheft zur hannoverschen Gradmessung aus dem Jahr 1825.

\par\noindent\hspace*{\parindent}Das Ausgleichungsverfahren von \textsc{Gauss} verläuft folgendermaßen.
Nachdem aus den Winkelgleichungen nach der Methode der kleinsten Quadrate
die Normalgleichungen auf S.~303 und 304 aufgestellt worden waren, löste
\textsc{Gauss} diese (auf indirektem Weg, vgl.~S.~325) und berechnete daraus
die Verbesserungen, die \textls[120]{allein} die Dreieckswinkelgleichungen erfüllen.
Mit den zum ersten Mal verbesserten Winkelwerten wurden die 12
Seitengleichungen in der für \textsc{Gauss} charakteristischen Weise
aufgestellt, S.~304---311. Für die Figur, auf die sich die Seitengleichung
bezieht, erreicht \textsc{Gauss} mit der modifizierten Seitengleichung, dass
die Ausgleichung, sofern diese Seitengleichung nicht mit einer anderen
zusammenhängt, stets nur die Winkelsummen der angrenzenden Dreiecke
beeinflusst und nicht die der inneren Dreiecke. Vgl.~den Brief an
\textsc{Gerling} vom 19.~Januar 1840, S.~253. Aus den umgeformten
Seitengleichungen wurden die 12 Normalgleichungen auf S.~312 hergeleitet.
Ihre Lösung liefert neue Werte für die Verbesserungen der Richtungen,
die \textls[120]{allein} die Seitengleichungen erfüllen. Nach dieser
ersten Ausgleichung werden die neu entstandenen Winkelschlussfehler der Dreiecke
wieder ausgeglichen. Mit den gefundenen Verbesserungen werden erneut die
Seitengleichungen und die zugehörigen Normalgleichungen aufgestellt, die
sich jetzt von den vorherigen nur in den konstanten Gliedern unterscheiden.
Ihre Lösung liefert wiederum neue Werte der Richtungsverbesserungen.
Auf diese Weise hat \textsc{Gauss} die Ausgleichung noch zweimal
wiederholt. Wie in Supplementum theor.~comb.~observ.~Art.~19 gezeigt
wird, führt dieses Verfahren zu denselben Werten wie die Ausgleichung
sämtlicher Winkel- und Seitengleichungen in einem einzigen Schritt. Dass
die Rechnung so schnell stabile Ergebnisse geliefert hat, liegt nach
\textsc{Gauss} an der von ihm gewählten Form der Seitengleichungen, vgl.~S.~251.

\par\noindent\hspace*{\parindent}Setzte man die Verbesserungen aus Art.~[7] in alle 51 Winkelgleichungen
ein, so ergäbe sich, wie aus der Zusammenstellung in Art.~[8] ersichtlich,
bei 11 Gleichungen der Fehler $0$, bei 18 Gleichungen der Fehler
$0^{\prime\prime}{,}001$, bei 18 Gleichungen der Fehler $0^{\prime\prime}{,}002$
und bei 4 Gleichungen der Fehler $0^{\prime\prime}{,}003$.

\par\noindent\hspace*{\parindent}\textsc{Gauss} hat die von den Punkten 28 bis 33 ausgehenden Richtungen
nochmals korrigiert. Woher diese Korrekturen stammen, ist nicht ersichtlich.
Wahrscheinlich gehen sie auf eine nachträgliche Änderung der beiden in
Langwarden beobachteten Winkel um $0^{\prime\prime}{,}128$ und
$0^{\prime\prime}{,}015$ zurück. Jedenfalls sind die unter [7] mitgeteilten
Verbesserungen, die zur Ableitung der folgenden Tabellen verwendet wurden,
diejenigen, die sich aus seiner Ausgleichung ergeben haben.

\par\noindent\hspace*{\parindent}Aus den 51 Winkelschlussfehlern der Dreiecke $w_i$ erhält man für den mittleren
Richtungsfehler $m$ nach der Näherungsformel
$\displaystyle m^2=\frac{1}{6\mathbin{.}51}\sum_{1}^{51}w_i^2$:
\[
m=\pm\,0^{\prime\prime}{,}643.
\]

\clearpage
\noindent328\hfill\textbf{BEMERKUNGEN. ZUR NETZAUSGLEICHUNG.}
\par\medskip

\par\noindent\hspace*{\parindent}Beim Anschluss des Punktes Hohenhorn an das Dreiecksnetz auf S.~316
mithilfe der Figur Hohenhorn.~21.~15.~17.~18 wurden die Seitenverhältnisse
und die Winkel der Gradmessung festgehalten. Die Dreiecksseite
Hamburg--Hohenhorn diente der hannoverschen Grad- und Landesvermessung
als Basis. Vgl.~Bestimmung des Breitenunterschiedes etc., S.~48, sowie
den Brief an \textsc{Schumacher} vom 22.~Dezember 1827, S.~281.

\par\noindent\hspace*{\parindent}Zu der Tabelle in Art.~[9] ist Folgendes zu bemerken. Die Darstellung
einer Dreiecksseite $i.k$ in der Ebene bilde mit der Geraden durch ihren
Anfangs- und Endpunkt die Winkel $\psi_{i.k}$ und $\psi_{k.i}$. Bezeichnet
$T$ das astronomische Azimut, $t$ das Azimut auf dem Sphäroid, $\Theta$
das Azimut in der Ebene und $\mathrm{c}$ die Meridiankonvergenz, so gilt
(vgl.~S.~201):
\[
\begin{array}{r@{\qquad}c}
1)&t_{i.k}=T_{i.k}+c_i\\[4pt]
2)&\Theta_{i.k}=t_{i.k}-\psi_{i.k},\qquad
\Theta_{k.i}=\Theta_{i.k}\pm180^\circ,\qquad
t_{k.i}=\Theta_{k.i}+\psi_{k.i}.
\end{array}
\]

\par\noindent\hspace*{\parindent}Die Berechnung von $\psi$ erfolgt nach der Formel in Art.~15, S.~159.
Vgl.~auch S.~216/217. Diese Formel setzt ebenso wie die zur Reduktion der
sphäroidischen Dreiecksseite auf die ebene die Kenntnis angenäherter
Koordinaten der Dreieckspunkte voraus, die man durch eine vorläufige
Berechnung mit geringer Genauigkeit bestimmen muss.

\par\noindent\hspace*{\parindent}Ist also im sphäroidischen Dreiecksnetz das astronomische Azimut
\textls[120]{einer} Dreiecksseite bekannt, so kann man durch wiederholte
Anwendung der Gleichungen (2) und mithilfe der Winkel des Netzes $t$ und
$\Theta$ für alle Seiten berechnen. Für alle Punkte auf dem Hauptmeridian
gilt nun $\mathrm{c}=0$. Geht daher die Dreiecksseite $1\,.\,2$, deren
Azimut beobachtet wurde, von einem Punkt dieses Meridians aus, so gilt
\[
\begin{array}{r@{\qquad}c}
3)&t_{1.2}=T_{1.2}.
\end{array}
\]

\par\noindent\hspace*{\parindent}Zur Orientierung des hannoverschen Dreieckssystems dient die Seite
Göttingen, Theodolithplatz 1823---Nördliches Meridianzeichen. Ihr Azimut
beträgt nach Art.~[9] $T_{1.2}=t_{1.2}=5^{\prime\prime}{,}473$.
In einem Brief an \textsc{Gerling} vom 26.~Dezember 1823
(vgl.~auch \textsc{Gerling}, Beiträge zur Geographie Kurhessens etc., S.~69)
gibt \textsc{Gauss} dafür $5^{\prime\prime}{,}471$ an. Aus den Koordinaten
des Theodolithplatzes von 1823 auf der Göttinger Sternwarte,
$x=-5{,}507\,\mathrm{m}$, $y=0$, und des nördlichen Meridianzeichens,
$x=-5019{,}756\,\mathrm{m}$, $y=-0{,}133\,\mathrm{m}$, Band~IV, S.~416/417,
ergibt sich ebenfalls das Azimut $=5^{\prime\prime}{,}471$.

\par\noindent\hspace*{\parindent}Sind nun die ebenen Azimute berechnet, wie in der Tabelle unter [9]
geschehen, so kann man aus ihnen die Winkel der ebenen Dreiecke
zusammenstellen und dann, wenn die lineare Länge einer Dreiecksseite
bekannt ist, die Längen aller übrigen Dreiecksseiten berechnen. Man muss
also nur \textls[120]{einmal} eine sphäroidische Dreiecksseite auf eine ebene
reduzieren. Diese Übertragung erfolgt nach der auf S.~215 angegebenen
Formel. Wendet man sie auf alle sphäroidischen Dreiecksseiten an, so
erhält man ebenfalls die ebenen Dreiecksseiten.

\par\noindent\hspace*{\parindent}Sind aber die Seiten der ebenen Dreiecke und ihre Azimute bekannt, so
lassen sich daraus leicht nacheinander die ebenen rechtwinkligen
Koordinaten der Dreieckspunkte berechnen.

\par\noindent\hspace*{\parindent}Die auf S.~299 wiedergegebene Karte des Dreiecksnetzes, das der
Ausgleichung unterzogen wurde, ist die Kopie einer Zeichnung, die
\textsc{Gauss} beim hannoverschen Kabinettsministerium
eingereicht hat. Im Original fehlen versehentlich die Richtungen
Wilsede-Brüttendorf und Timpenberg-Hamburg. Es enthält dagegen noch
den Anschluss der Punkte Hohenhorn und Lauenburg sowie den der Punkte
Wangeroog und Neuwerk.

\par\noindent\hspace*{\parindent}Zur Ausgleichung der in den ersten Jahren der Gradmessung beobachteten
Dreiecke vgl.~die weiter unten folgenden Briefe an \textsc{Bessel} vom
5.~November 1823 und an \textsc{Bohnenberger} vom 16.~November 1823.
\hfill\textsc{Krüger.}

\begin{center}\rule{.24\textwidth}{.4pt}\end{center}
